\documentclass[11pt,a4paper]{article}
\usepackage[T1]{fontenc}
\usepackage[utf8]{inputenc}
\usepackage[main=italian,provide=*]{babel}
\usepackage{amsmath,amssymb}
\usepackage{geometry}
\geometry{margin=2.3cm,top=2.5cm,bottom=2.7cm}
\usepackage{booktabs}
\usepackage[table]{xcolor}
\usepackage{tcolorbox}
\tcbuselibrary{skins,breakable}
\usepackage{titlesec}
\usepackage{enumitem}
\usepackage{needspace}
\usepackage{fancyhdr}
\usepackage{hyperref}
\hypersetup{colorlinks=true,linkcolor=Sepia,citecolor=Sepia,urlcolor=Sepia}
\newcommand{\Tr}{\operatorname{Tr}}
\newcommand{\ket}[1]{\lvert #1 \rangle}
\newcommand{\bra}[1]{\langle #1 \rvert}

% --- palette ---
\definecolor{inkblue}{HTML}{1B3A5C}
\definecolor{lightblue}{HTML}{EAF1F8}
\definecolor{accent}{HTML}{B0562A}
\definecolor{softgray}{HTML}{6B6B6B}
\definecolor{okgreen}{HTML}{2E6B3E}
\definecolor{okgreenbg}{HTML}{EAF5EC}

% --- titoli di sezione ---
\titleformat{\section}
  {\normalfont\Large\bfseries\color{inkblue}}
  {\thesection}{0.6em}{}[{\color{inkblue!40}\titlerule}]
\titlespacing{\section}{0pt}{0.7em}{0.4em}
\titleformat{\subsection}
  {\normalfont\large\bfseries\color{accent}}
  {}{0em}{}
\titlespacing{\subsection}{0pt}{0.5em}{0.2em}

\pagestyle{fancy}
\fancyhf{}
\renewcommand{\headrulewidth}{0.4pt}
\fancyhead[L]{\small\color{softgray}Scan sui parametri di rumore --- Passo 5, Parte 2}
\fancyhead[R]{\small\color{softgray}\thepage}
\fancyfoot[C]{\small\color{softgray}Tesi triennale, Samuele --- Relatore: Prof.\ Alessandro Chiesa}

% --- box riutilizzabili ---
\newtcolorbox{keyeq}[1][]{
  colback=lightblue, colframe=inkblue!70, boxrule=0.6pt,
  arc=2pt, left=8pt, right=8pt, top=3pt, bottom=3pt, #1
}
\newtcolorbox{okbox}[1]{
  colback=okgreenbg, colframe=okgreen!60, boxrule=0.6pt,
  arc=2pt, left=7pt, right=7pt, top=2pt, bottom=2pt,
  fonttitle=\bfseries\color{okgreen}, title={#1}
}

\title{\vspace{-1.5em}\color{inkblue}Scan sui parametri di rumore --- Passo 5 (ultimo)\\[0.2em]
\Large\normalfont\color{softgray}Come si sposta $N^*$: chiusura della Parte 2}
\author{Note di lavoro --- Tesi triennale, Samuele\\ Relatore: Prof.\ Alessandro Chiesa}
\date{}

\begin{document}
\sloppy
\maketitle
\thispagestyle{fancy}
\vspace{-2.2em}

File: \texttt{scan\_parametri\_rumore\_dimero.py},
\texttt{validate\_scan\_parametri\_rumore\_dimero.py}. Osservabile: fedeltà
di Trotter (Passo 3) per gli scan su $\varepsilon_{1q},\varepsilon_{2q}$;
correlatore (Passo 4) per la verifica su $p_\text{readout}$.

\section*{Un risultato analitico prima dello scan: $p_\text{readout}$ non sposta $N^*$}
\addcontentsline{toc}{section}{Risultato analitico su $p_\text{readout}$}

La correzione di readout derivata al Passo 4, $\langle Z\rangle_\text{readout}=
(1-2p)\langle Z\rangle_\text{ideale}$, è un fattore \emph{moltiplicativo
costante in $N$}. Se una funzione $g(N)$ ha un massimo a $N^*$, allora
$c\cdot g(N)$ con $c>0$ costante ha il massimo esattamente allo stesso
$N^*$ --- moltiplicare per una costante positiva non può spostare un
argmax. Quindi:
\begin{keyeq}
\centering
$N^*$ non dipende da $p_\text{readout}$ (per un readout simmetrico)
\end{keyeq}
Verificato numericamente su $C_{21}^{xx}(t=2)$: $N^*=5$ per
$p_\text{readout}=0, 0.05, 0.20$ (fattore 4 di variazione), con i valori
del correlatore che scalano proporzionalmente ma il massimo sempre nello
stesso punto. Questo riduce lo scan effettivo a due soli parametri:
$\varepsilon_{1q}$ ed $\varepsilon_{2q}$.

\section*{Scan su $\varepsilon_{2q}$}
\addcontentsline{toc}{section}{Scan su $\varepsilon_{2q}$}

A parità di $\varepsilon_{1q}=2.9\times10^{-4}$ (riferimento):

\begin{center}
\renewcommand{\arraystretch}{1.2}
\begin{tabular}{@{}ccc@{}}
\toprule
$\varepsilon_{2q}$ & $N^*$ & $F(N^*)$ \\
\midrule
$1\times10^{-3}$ & $10$ & $0.879861$ \\
$2\times10^{-3}$ & $9$ & $0.855726$ \\
$3.8\times10^{-3}$ (riferimento) & $8$ & $0.817725$ \\
$6\times10^{-3}$ & $7$ & $0.777595$ \\
$1\times10^{-2}$ & $1$ & $0.736025$ \\
$1.5\times10^{-2}$ & $1$ & $0.719825$ \\
$2.5\times10^{-2}$ & $1$ & $0.688722$ \\
\bottomrule
\end{tabular}
\end{center}

\section*{Scan su $\varepsilon_{1q}$}
\addcontentsline{toc}{section}{Scan su $\varepsilon_{1q}$}

A parità di $\varepsilon_{2q}=3.8\times10^{-3}$ (riferimento). Il singolo
passo ha 3 CNOT ma anche 15 gate a 1 qubit (8 $R_z$ + 7 $\sqrt X$): il
contributo di $\varepsilon_{1q}$ non è trascurabile a priori, va misurato:

\begin{center}
\renewcommand{\arraystretch}{1.2}
\begin{tabular}{@{}ccc@{}}
\toprule
$\varepsilon_{1q}$ & $N^*$ & $F(N^*)$ \\
\midrule
$1\times10^{-4}$ & $9$ & $0.844729$ \\
$2.9\times10^{-4}$ (riferimento) & $8$ & $0.817725$ \\
$1\times10^{-3}$ & $7$ & $0.733734$ \\
$3\times10^{-3}$ & $1$ & $0.672574$ \\
$1\times10^{-2}$ & $1$ & $0.513302$ \\
$3\times10^{-2}$ & $1$ & $0.315991$ \\
\bottomrule
\end{tabular}
\end{center}

\begin{okbox}{Tre controlli superati}
\begin{itemize}[leftmargin=1.2em,itemsep=0.2em]
\item \textbf{Nessuna regressione}: $N^*$ al punto di riferimento
ricalcolato con la funzione generale di questo modulo coincide con quello
trovato manualmente al Passo 3 ($N^*=8$).
\item \textbf{Monotonia}: $N^*$ è non-crescente su tutto lo scan, sia in
$\varepsilon_{1q}$ sia in $\varepsilon_{2q}$ --- più rumore, meno passi
convenienti, mai il contrario.
\item \textbf{Invarianza da $p_\text{readout}$}: dimostrata sopra,
confermata numericamente.
\end{itemize}
\end{okbox}

\needspace{8\baselineskip}
\section*{Osservazione: saturazione a $N^*=1$}
\addcontentsline{toc}{section}{Osservazione}

Oltre una certa soglia di rumore, $N^*$ satura a $1$: il compromesso
smette di avere un ottimo intermedio, e conviene il passo più grezzo
possibile piuttosto che pagare il costo di CNOT aggiuntivi. Da notare
onestamente: il modello ingenuo $\varepsilon_\text{tot}(N)\sim a/N+
g\varepsilon_{2q}N$ prevede $N^*\propto1/\sqrt{\varepsilon_{2q}}$; lo
spostamento osservato è più lento di questa legge (un fit sui punti non
saturati, $\varepsilon_{2q}=10^{-3}$--$6\times10^{-3}$, dà una pendenza
log-log $\approx-0.2$, non $-0.5$). Non è una contraddizione: quel modello
era pensato per un errore additivo su una singola quantità, mentre qui
l'osservabile è una fedeltà, che con rumore depolarizzante accumulato su
molti passi si comporta più come un decadimento (moltiplicativo) che come
una somma --- la forma funzionale esatta di $N^*(\varepsilon)$ per la
fedeltà resta un affinamento possibile, non necessario per rispondere alla
domanda del relatore (``come si sposta $N^*$''), a cui questa tabella
risponde direttamente.

\needspace{6\baselineskip}
\section*{Parte 2: chiusura}
\addcontentsline{toc}{section}{Parte 2: chiusura}

Con questo si chiudono tutti e cinque i passi della pipeline di Parte 2
sul dimero (modello di rumore, VQE+DM rumoroso, Trotter rumoroso,
correlatori rumorosi, scan sui parametri). Prossimo passo: integrazione di
tutti i risultati nel documento finale di tesi, incluso valorizzare il
confronto RBS vs $W$ come sezione metodologica a sé (nota già presente in
\texttt{scheda\_progetto\_tesi.md}).

\end{document}
